\documentclass[conference]{IEEEtran}

% ============================================================
% PACKAGES
% ============================================================
\usepackage[hidelinks]{hyperref}
\usepackage{cite}
\usepackage{amsmath,amssymb,amsfonts}
\usepackage{algorithmic}
\usepackage{algorithm}
\usepackage{array}
\usepackage{booktabs}
\usepackage{multirow}
\usepackage{url}
\usepackage{graphicx}
\usepackage{xcolor}
\usepackage{balance}
\usepackage{float}

% ============================================================
% TITLE
% ============================================================

\title{
Automated Difficulty and Algorithmic Tag Prediction for
Competitive Programming Problems Using Natural Language Processing
}

\author{
\vspace{4cm}

}
% ============================================================
% DOCUMENT
% ============================================================

\begin{document}

\maketitle

% ============================================================
% ABSTRACT
% ============================================================
\textbf{Abstract}  Competitive programming platforms contain huge pools of algorithmic problems
with difficulty levels and algorithmic tags for problem selection, personalised
practice, and intelligent education. This paper proposes a lightweight Natural
Language Processing (NLP) framework that predicts the difficulty rating and
algorithmic tags for Codeforces problems directly from their textual
descriptions.
The proposed system combines the problem title, statement, input description,
output description, and available tag information into a single textual
representation. The text data is converted into numerical feature vectors with
up to 20,000 unigram and bigram features using Term Frequency--Inverse Document
Frequency (TF--IDF). For difficulty prediction, a Random Forest regression
model is trained to estimate the problem rating.
The dataset contains 7,185 Codeforces problems, with 5,000 used for training
and 100 held-out problems used for evaluation.
The model predicted 73 out of 100 problems within $\pm100$ rating points and
81 out of 100 problems within $\pm200$ rating points, corresponding to
$73.00\%$ and $81.00\%$, respectively.
For this test set, the Random Forest model produced a Mean Absolute Error
(MAE) of 112.35 and a Root Mean Squared Error (RMSE) of 254.01. The tag
prediction experiment, evaluated on the same 100 held-out problems, achieved
$2.50\%$ exact-match accuracy, $82.19\%$ at-least-one-correct-tag accuracy,
$32.28\%$ micro-precision, $50.00\%$ micro-recall, and $39.70\%$ micro-F1
score.
These results suggest that lightweight lexical representations can help predict
competitive programming metadata, but also emphasise the difficulty of
recovering abstract algorithmic concepts from natural language descriptions
alone.

% ============================================================
% KEYWORDS
% ============================================================

\begin{IEEEkeywords}
Competitive Programming, Codeforces, 
TF--IDF, Random Forest, Difficulty Prediction, Algorithmic Tag Prediction.

\end{IEEEkeywords}

% ============================================================
% ============================================================
% INTRODUCTION
% ============================================================

\section{\underline{Introduction}}

Competitive programming plays a significant role in developing algorithmic thinking, mathematical reasoning, problem solving skills and coding implementation abilities. Online platforms like Codeforces host thousands of programming problems covering large scale of topics, such as, dynamic programming, graph algorithm, greedy methods, number theory and data structures \cite{codeforcesapi}.

Each problem on this platform is associated with specific metadata notably a difficulty rating and algorithmic tags \cite{codeforcesapi}.

The difficulty ratings are especially helpful to the learners because they let the programmers choose problems that match their current skill level. Algorithmic tags are also useful in that they can enable users to practise particular concepts and they can allow systems to organise problems by topic.

However for a newly published problem this metadata is not immediately available. It is usually determined after many users solve it. If the difficulty level and algorithmic tag could be automatically predicted from the problem text description it would be beneficial for educational platforms and recommendation system.

A competitive programming problem transforms a computational challenge
into a structured statement that can be understood, analyzed, and solved
algorithmically.A competitive programming problem normally contains several textual
components, including a title, problem statement, input specification,
output specification, constraints, and examples. These components provide
different types of information about the computational task.For example, the statement may indicate that a solution must find a
minimum distance, while the constraints may indicate that the input size
can reach $10^5$. Together, these textual signals can provide information
about the likely algorithmic complexity of the required solution.

\begin{figure}[htbp]
\centering
\includegraphics[width=0.85\linewidth]{cf.png}
\caption{Structure and textual components of a Codeforces programming problem.}
\label{fig:cp_problem_structure}
\end{figure}

The Codeforces platform associates programming problems with numerical
difficulty ratings \cite{codeforcesapi}. In this work, difficulty estimation
is treated as a regression task.

A programming problem may require multiple algorithmic concepts. For
example, a problem may simultaneously involve graph algorithms and
shortest paths. Therefore, tag prediction is formulated as a multi-label
classification problem \cite{tsoumakas2007}.

% In these paper we analyzed whether classical or traditional NLP and machine learning techniques can extract from the text of competitive programming problem.

For this we proposed a lightweight and computationally efficient pipeline based on TF-IDF text representation \cite{salton1988,manning2008} and a Random Forest Regressor \cite{breiman2001}.

% 3 LITERATURE REVIEW
% ============================================================

\section{\underline{Literature Review}}

Statistical text representations remain useful when computational cost and
reproducibility are important. Term Frequency--Inverse Document Frequency
(TF--IDF) represents a document through terms that are frequent locally but
comparatively rare across the collection \cite{salton1988,manning2008}.
McCallum and Nigam also showed that comparatively simple probabilistic models
can provide strong bag-of-words baselines for text classification
\cite{mccallum1998}. These methods do not explicitly represent program
semantics, but they provide an efficient reference point against which more
complex neural systems can be evaluated.

\subsection{Machine Learning for Regression and Classification}

For predictive modelling, Random Forests aggregate many decision trees and
can capture non-linear feature interactions \cite{breiman2001}. Logistic
regression is also a well-established linear classifier for sparse textual
features. In a one-vs-rest formulation, one binary classifier is learned for
each label, making it a natural baseline for multi-label text categorisation
\cite{tsoumakas2007}. These properties motivate the two classical models used
in CPVision-AI: Random Forest regression for numerical ratings and one-vs-rest
logistic regression for algorithmic tags.

\subsection{Algorithmic Classification from Problem Statements}

Research directly concerned with programming problems has established
algorithm-class prediction as a text-classification task. Athavale et al.
introduced four datasets derived from programming word problems, including
multi-class and multi-label Codeforces datasets, and compared neural and
non-neural classifiers \cite{athavale2019}. Their experiments demonstrate that
problem statements contain useful algorithmic signals, while their comparison
with human predictions also illustrates the ambiguity of inferring a solution
method from text alone.

Iancu et al. studied automatic multi-label tagging of programming challenges
from Codeforces and Topcoder using TF--IDF, Doc2Vec, and recurrent neural
representations \cite{iancu2019}. Their neural models generally improved upon
the information-retrieval baselines, suggesting that semantic representations
can benefit this task. More recently, Lokat et al. evaluated MLP, LSTM, and GRU
models on 2,400 Codeforces statements and reported the MLP as their strongest
model \cite{lokat2023}. Together, these studies support both the feasibility of
text-based tag prediction and the value of retaining inexpensive lexical
baselines.

\subsection{Difficulty Estimation and Educational Applications}

Difficulty in an online judge can be inferred from content, learner behaviour,
or a combination of the two. Zhao et al. jointly investigated topic discovery
and difficulty estimation from large-scale interaction traces in online judge
systems \cite{zhao2018}. Their competition-based expertise model used learned
topics and user attempts, demonstrating the value of behavioural evidence when
historical interaction data are available. CPVision-AI instead addresses the
cold-start setting in which a newly added problem has textual content but no
submission history.

Yoshida et al. recently examined difficulty estimation for programming
problems using a fine-tuned large language model \cite{yoshida2024}. This work
represents the shift from hand-designed or sparse features toward pretrained
semantic representations for rating prediction. In a related educational
setting, Xu et al. represented programming exercises through difficulty,
knowledge concepts, problem text, input, and output, using a bidirectional GRU
within a reinforcement-learning recommendation framework \cite{xu2022}. These
studies show that difficulty and topic metadata are useful not only for archive
organisation but also for adaptive practice and personalised recommendation.

\subsection{Deep Learning and Source Code Analysis}

Recent NLP research has shifted toward contextual transformer representations.
BERT uses bidirectional pre-training to model contextual relationships in
natural language \cite{devlin2019}, whereas CodeBERT extends pre-training to
paired natural-language and programming-language data \cite{feng2020}.
Lobanov et al. applied this distinction directly to competitive-programming tag
prediction by comparing statement-based and solution-code-based approaches and
then combining BERT statement features with a graph neural representation of
source code \cite{lobanov2023}. Their ensemble outperformed approaches using
only one information source, showing that submitted solutions provide semantic
evidence that may not be recoverable from a statement alone. Such models can
improve representation quality, but training and deploying them generally
requires substantially more computation than sparse lexical models
\cite{strubell2019}.

\subsection{Automated Problem Analysis in Competitive Programming}

The work closest to the present study is the Problem-Solving Guide of Kim
et al., which introduced the AMT dataset and a deep multi-task architecture for
jointly predicting algorithm tags and difficulty from competitive-programming
problems \cite{kim2023}. In contrast, CPVision-AI studies the same two outputs
with independently trained classical models and deploys them in a lightweight
web application. This design provides an inexpensive and reproducible baseline
that uses only information available in the title, statement, input, and output
descriptions at inference time. Unlike approaches that analyse accepted
solutions \cite{lobanov2023}, it is therefore applicable before a solution has
been submitted. Its principal limitation is that TF--IDF captures lexical
correlation rather than the deeper reasoning required to derive an algorithm.

% The numerical difficulty distribution is summarized in
% Table~\ref{tab:dataset}.

% All 7,185 records in the prepared dataset contain a numerical difficulty
% rating. For the main difficulty prediction experiment, 5,000 problems were
% selected for model development and 100 problems were independently selected
% as unseen test.
\begin{figure}[htbp]
    \centering
    \includegraphics[width=0.95\linewidth]{Rating.png}
    \caption{Distribution of difficulty ratings across the 7,185 Codeforces problems.}
    \label{fig:rating_distribution}
\end{figure}
% ============================================================
% 5 METHODOLOGY
% ============================================================

\section{\underline{Methodology}}

The proposed methodology combines text preparation, task-specific TF--IDF
feature extraction, multi-label algorithmic-tag classification, and
difficulty-rating regression. The stages and their relationships are
summarized in Fig.~\ref{fig:cpvision_framework} and described in the
following subsections.

\subsection{Proposed CPVision-AI Framework}

\begin{figure}[H]
    \centering
    \includegraphics[width=0.98\columnwidth]{flow.png}
    \caption{Overall CPVision-AI framework for difficulty-rating estimation
    and multi-label algorithmic-tag prediction.}
    \label{fig:cpvision_framework}
\end{figure}

CPVision-AI receives the title, statement, input description, and output
description of a Codeforces problem. After the textual fields are cleaned
and combined, two task-specific TF--IDF representations are produced. The
first is processed by a one-vs-rest logistic regression classifier to
predict multiple algorithmic tags. The second is processed by a Random
Forest regressor to estimate the difficulty rating and corresponding rank.
During joint inference, the predicted tags provide additional evidence to
the difficulty-prediction branch.

% \begin{figure}[H]
%     \centering
%     \includegraphics[width=0.98\columnwidth]{flow.png}
%     \caption{Overall CPVision-AI framework for difficulty-rating estimation
%     and multi-label algorithmic-tag prediction.}
%     \label{fig:cpvision_framework}
% \end{figure}

\subsection{Dataset Collection}

The dataset contains 7,185 Codeforces programming problems and was derived
from the public \emph{Codeforces Competitive Programming Dataset} hosted on
Kaggle \cite{kagglecodeforces}. Codeforces is a popular competitive
programming platform where programmers improve their problem-solving,
logical thinking, and debugging skills by participating in programming
contests \cite{codeforcesapi}.

Each contest contains problems with different difficulty levels and algorithmic concepts. Each record in the dataset contains several textual and metadata fields, including the problem ID, index, title, statement, input description, output description, algorithmic tags, and numerical difficulty rating.

The main attributes used by the proposed system are:

\begin{itemize}
\item problem title,
\item problem statement,
\item input description,
\item output description,
\item problem tags, and
\item numerical difficulty rating.
\end{itemize}

These attributes were used as the primary source of information for developing the CPVision-AI difficulty prediction and algorithmic tag prediction models. The textual fields were combined to form a unified representation of each problem, allowing the models to learn patterns associated with problem difficulty and algorithmic categories.


\subsection{Dataset Pre-processing}
The original Kaggle dataset \cite{kagglecodeforces} contained several fields
that were not required by CPVision-AI. Therefore, the dataset was customised
for the two prediction tasks. During preprocessing, irrelevant fields were
removed and missing or incomplete records were handled.

For our final dataset we kept the problem ID, problem name, problem index, problem statement, tags, rating, editorial, input description, output description and solution. These fields were found to be useful for difficulty prediction, algorithmic tag prediction and related analysis of competitive programming problems.

After preprocessing and customisation, the final dataset comprised 7,185
Codeforces problems with difficulty ratings ranging from 800 to 3{,}500, each
associated with one or more algorithmic tags. For model development, 5,000
problems were used for training; 100 separate problems not seen during
training were held out as the evaluation set. The remaining 2,085 problems
were excluded from both training and evaluation to maintain a clean
separation between development and test data.

\subsection{TF-IDF feature extraction}

For transforming the combined textual representation of each problem into numerical feature vectors, we used Term Frequency--Inverse Document Frequency (TF--IDF) \cite{salton1988,manning2008}. TF--IDF gives more weight to words that appear frequently in a given problem, but are not that common in the whole dataset.

Term Frequency--Inverse Document Frequency (TF--IDF) is used to convert
the combined textual representation into numerical feature vectors.
\begin{figure}[htbp]
    \centering
    \includegraphics[width=0.85\linewidth]{statements.png}
    \caption{Distribution of problem-statement lengths in the Codeforces dataset.}
    \label{fig:statement_length_distribution}
\end{figure}
For a term $t$ and document $d$, the TF--IDF value is calculated as:
\begin{equation}
TFIDF(t,d) =
TF(t,d)
\left(
\log
\left(
\frac{N+1}{DF(t)+1}
\right)
+1
\right)
\end{equation}

where $N$ is the number of documents and $DF(t)$ is the number of documents
containing term $t$ \cite{salton1988}.

The implementation uses:

\begin{itemize}
    \item maximum 20,000 features,
    \item unigram and bigram features,
    \item minimum document frequency of 2,
    \item maximum document frequency of 0.95, and
\end{itemize}
\vspace{\baselineskip}

We used Bigrams as it is important because many programming concepts are expressed using multiple words, such as “Binary search”, “Shortest path”, “Dynamic Programming”, and “Brute force”.

\subsection{Random Forest Regression for Rating Prediction}

Random Forest Regression was used to predict the difficulty rating of
Codeforces problems. Random Forest is an ensemble learning method that
combines the predictions of multiple decision trees instead of relying on
a single tree \cite{breiman2001}. In our system, the textual information of each problem is first converted
into numerical features using TF--IDF. These features are then given to the
Random Forest Regressor, which learns the relationship between the textual
features and the Codeforces difficulty rating. Random Forest was selected
because the relationship between problem text and difficulty may not always
be linear.

Let $x_i$ denote the TF--IDF feature vector of problem $i$, and let $y_i$
denote its numerical difficulty rating. The regression objective is to learn
a function $f$ that produces the predicted rating $\hat{y}_i$:

\begin{equation}
f(x_i) \rightarrow \hat{y}_i.
\end{equation}

For an input feature vector $x$, the predicted rating is:

\begin{equation}
\hat{y}
=
\frac{1}{B}
\sum_{b=1}^{B}
T_b(x)
\end{equation}

where $B$ is the number of trees and $T_b(x)$ represents the prediction of
the $b$-th decision tree \cite{breiman2001}.

Random Forest was also suitable for our task because it can work with a
large number of textual features and learn different patterns from the
problem statements. The implementation was developed using scikit-learn
\cite{sklearn2011}.

The model hyperparameters were selected through 5-fold cross-validation on
the 5,000 training problems. The configuration that minimised validation MAE
consisted of 150 decision trees with no constraint on maximum tree depth,
square-root feature selection at each split, and a fixed random seed of 42
for reproducibility. Bootstrapping was disabled because it did not improve
validation performance on this dataset.

\subsection{Multi-Label Tag Prediction}
This step entails the task of predicting the algorithmic tags of a competitive programming problem. This is because the algorithms used in solving a certain problem can vary and therefore a single problem may be associated with several algorithmic tags at once \cite{tsoumakas2007}. For instance, a problem can have tags such as “graphs”, “greedy”, and “shortest paths” at once.

Let $\mathcal{T}$ denote the complete set of supported algorithmic tags.
For problem $i$, the ground-truth tag set $Y_i$ satisfies

\begin{equation}
Y_i \subseteq \mathcal{T}.
\end{equation}

Given the TF--IDF feature vector $x_i$, the multi-label objective is to
predict the tag set $\hat{Y}_i$ using a function $g$:

\begin{equation}
g(x_i) \rightarrow \hat{Y}_i.
\end{equation}

The representation of the problem statements numerically is done by firstly transforming the features using the term frequency-inverse document frequency (TF-IDF) \cite{salton1988,manning2008}. The vectors obtained are then fed to the one vs. rest (OvR) logistic regression classifier implemented using scikit-learn \cite{sklearn2011}. Instead of having one classifier to classify the tags, we train different binary classifiers for each of the algorithmic tags that is being predicted. Each classifier determines whether the respective tag is applicable to the problem or not.
\begin{figure}[htbp]
    \centering
    \includegraphics[width=0.95\linewidth]{tags.png}
    \caption{Frequency distribution of algorithmic tags in the Codeforces dataset.}
    \label{fig:tag_distribution}
\end{figure}

Since the classifiers work independently, several classifiers can give positive classifications to a single problem which makes the model capable of classifying several algorithmic tags to a single problem.


% ============================================================
% 6 EXPERIMENTAL SETUP
% ============================================================

% \section{\underline{Experimental Setup}}

% \subsection{Train--Test Separation}

% A major objective of the experiment is to evaluate the ability of the model
% to generalize to problems that were not used during training.

% For the main difficulty experiment, 5,000 problems were selected for model
% training. A separate set of 100 problems was then selected from problems
% that were not part of the training set.

% The TF--IDF vectorizer was fitted on the training text only. The resulting
% vectorizer was subsequently used to transform the 100 unseen problems.

% This procedure prevents the test documents from influencing the learned
% vocabulary or feature statistics.

% The experimental pipeline can therefore be summarized as:

% \begin{equation}
% \text{Training Text}
% \rightarrow
% \text{TF--IDF Fit}
% \rightarrow
% \text{Random Forest Training}
% \end{equation}

% followed by:

% \begin{equation}
% \text{Unseen Text}
% \rightarrow
% \text{Existing TF--IDF}
% \rightarrow
% \text{Random Forest Prediction}
% \end{equation}

% \subsection{Regression Evaluation Metrics}

% The difficulty prediction model is evaluated using Mean Absolute Error
% (MAE):

% \begin{equation}
% MAE =
% \frac{1}{n}
% \sum_{i=1}^{n}
% |y_i-\hat{y}_i|
% \end{equation}

% Root Mean Squared Error (RMSE) is calculated as:

% \begin{equation}
% RMSE =
% \sqrt{
% \frac{1}{n}
% \sum_{i=1}^{n}
% (y_i-\hat{y}_i)^2
% }
% \end{equation}

% Because competitive programming ratings are commonly interpreted in
% rating intervals, tolerance-based accuracy is also reported:

% \begin{equation}
% Accuracy_{\pm k}
% =
% \frac{1}{n}
% \sum_{i=1}^{n}
% \mathbb{I}
% (|y_i-\hat{y}_i|\leq k)
% \times 100
% \end{equation}

% where $k$ is either 100 or 200 rating points.

% \subsection{Multi-Label Evaluation}

% For tag prediction, the following metrics are used.

% Micro-precision is:

% \begin{equation}
% P_{\mu}
% =
% \frac{
% \sum_j TP_j
% }{
% \sum_j(TP_j+FP_j)
% }
% \end{equation}

% Micro-recall is:

% \begin{equation}
% R_{\mu}
% =
% \frac{
% \sum_j TP_j
% }{
% \sum_j(TP_j+FN_j)
% }
% \end{equation}

% Micro-F1 is:

% \begin{equation}
% F1_{\mu}
% =
% 2
% \frac{
% P_{\mu}R_{\mu}
% }{
% P_{\mu}+R_{\mu}
% }
% \end{equation}

% Exact-match accuracy requires the complete predicted tag set to match the
% ground-truth tag set.

% ============================================================
% 7 RESULTS
% ============================================================
\section{\underline{Results and Discussion}}

\subsection{Difficulty Rating Prediction}

The Random Forest model was tested on 100 problems that were not used during
training. The results are shown in Table~\ref{tab:difficulty_metrics}.

\begin{table}[htbp]
\caption{Difficulty Prediction on 100 Unseen Problems}
\centering
\begin{tabular}{lcc}
\toprule
Metric & Random Forest & Mean Baseline \\
\midrule
MAE & 112.35 & 285.60 \\
RMSE & 254.01 & 389.74 \\
Accuracy within $\pm100$ & 73.00\% & 18.00\% \\
Accuracy within $\pm200$ & 81.00\% & 34.00\% \\
\bottomrule
\end{tabular}
\label{tab:difficulty_metrics}
\end{table}

The Mean Baseline always predicts the mean training rating; it serves as a
trivial lower bound. Compared with this baseline, the Random Forest
reduces MAE by 60.5\% (from 285.60 to 112.35) and increases tolerance
accuracy within $\pm100$ points from 18\% to 73\%, demonstrating that
textual features provide substantial predictive signal beyond the global
mean.

The RMSE of 254.01 indicates that some problems had larger prediction
errors. Words such as ``subarray'', ``permutation'', ``minimum distance'',
and ``binary search'', as well as input constraints, can give useful hints
about the type and difficulty of a problem. However, the model does not
understand the underlying solution algorithm, so certain problems remain
difficult to predict correctly.

\subsection{Example Difficulty Predictions}

Table~\ref{tab:example_predictions} shows some prediction examples.

\begin{table}[htbp]
\caption{Example Difficulty Predictions}
\centering
\begin{tabular}{lccc}
\toprule
Problem & Actual & Predicted & Error \\
\midrule
G. Vasya and Maximum Profit & 2400 & 2400 & 0 \\
B. Shuffle & 1300 & 1303 & 3 \\
B. Accordion & 1300 & 1287 & 13 \\
D. Dogeforces & 2300 & 2269 & 31 \\
E. Sleeping Schedule & 1700 & 1609 & 91 \\
\bottomrule
\end{tabular}
\label{tab:example_predictions}
\end{table}

These examples show that the model can predict some ratings very closely.
However, the overall RMSE also shows that larger errors occurred for some
other problems.

\subsection{Algorithmic Tag Prediction}

The tag prediction model was evaluated on the same 100 held-out problems
used for difficulty prediction. The results are shown in
Table~\ref{tab:tag_metrics}.

\begin{table}[htbp]
\caption{Multi-Label Algorithmic Tag Prediction (100 Held-Out Problems)}
\centering
\begin{tabular}{lcc}
\toprule
Metric & OvR LR & Majority Baseline \\
\midrule
Exact Match Accuracy & 2.50\% & 0.00\% \\
Micro-Precision & 32.28\% & 11.40\% \\
Micro-Recall & 50.00\% & 100.00\% \\
Micro-F1 & 39.70\% & 20.42\% \\
Hamming Loss & 0.1433 & 0.3800 \\
At-Least-One-Correct-Tag Accuracy & 82.19\% & 100.00\% \\
\bottomrule
\end{tabular}
\label{tab:tag_metrics}
\end{table}

The Majority Baseline always predicts the top-$k$ most frequent tags (where
$k$ equals the average ground-truth tag count). It achieves perfect recall
but very low precision, giving a micro-F1 of 20.42\%. The OvR logistic
regression classifier nearly doubles the micro-F1 to 39.70\% by trading
some recall for substantially higher precision (32.28\% vs.\ 11.40\%),
confirming that lexical features provide genuine discriminative signal.

The Hamming loss of 0.1433 indicates that 14.33\% of individual
label decisions were incorrect. Exact-match accuracy was only 2.50\%,
because this strict metric requires the complete predicted tag set to equal
the ground-truth set. Nevertheless, at least one correct tag was predicted
for 82.19\% of the test problems, demonstrating that partial tag
identification is considerably easier than recovering every applicable tag
without adding any incorrect labels.

Tag prediction is inherently difficult because the same words can appear in
problems that require different algorithms. TF--IDF captures lexical
co-occurrence but cannot model the deeper reasoning required to determine
a solution algorithm.

\subsection{Unseen Evaluation and Limitations of the Test Set}
Both models were trained on 5,000 problems and evaluated on 100 separate
problems not seen during training. The remaining 2,085 problems were
excluded from both splits to maintain a strict development/test boundary.

Although using held-out problems provides a direct measure of
generalisation, a test set of 100 problems represents only 1.4\% of the
full dataset. At this size, individual outlier problems can noticeably shift
the reported MAE and F1 values, limiting the statistical reliability of the
results. The dataset contains 7,185 problems, so a larger held-out set of
500--1{,}000 problems would substantially improve estimate stability at
minimal additional cost. Extending the test set and adopting
cross-validation are the primary recommendations for future work.



% ============================================================
% 9 DEPLOYMENT
% ============================================================

\section{\underline{Model Deployment and Inference}}

The trained difficulty prediction model and TF--IDF vectorizer can be saved
and reused without retraining. The machine-learning components are implemented using scikit-learn \cite{sklearn2011}.

The saved components consist of:

\begin{itemize}
    \item the trained Random Forest rating model, and
    \item the fitted TF--IDF vectorizer.
\end{itemize}

For a new problem, the system accepts the problem statement as input,
transforms it using the previously fitted TF--IDF vectorizer, and passes
the resulting feature vector to the Random Forest model.

\begin{figure}[htbp]
    \centering
    \includegraphics[width=0.85\linewidth]{web.png}
    \caption{CPVision-AI web interface showing predicted difficulty and algorithmic tags.}
    \label{fig:web_prediction_interface}
\end{figure}

This makes the trained model suitable for integration into a larger
CPVision-AI application.

An important requirement is that the same TF--IDF vectorizer used during
training must be used during inference. A new vectorizer must not be fitted
on the individual test or user-provided problem.

% ============================================================
% 10 LIMITATIONS
% ============================================================

% \section{\underline{Limitations}}

% Several limitations should be considered when interpreting the results.

% First, the difficulty experiment uses a test set of only 100 unseen
% problems. Although this provides a direct evaluation of generalization,
% larger test sets would produce more stable estimates.

% Second, TF--IDF primarily captures lexical information. It does not
% explicitly model the logical structure of a problem or understand the
% algorithmic solution required to solve it.

% Third, Random Forest regression is not specifically designed for high-
% dimensional natural-language representations. Linear models and
% transformer-based language models would provide useful comparison
% baselines in future experiments.

% Fourth, the algorithmic tag task is inherently multi-label and highly
% variable. Some tags occur much more frequently than others, which can make
% rare-tag prediction difficult.

% Finally, numerical difficulty ratings are not necessarily determined
% entirely by textual characteristics. Factors such as implementation
% complexity, contest context, mathematical insight, and participant
% performance can also influence the assigned rating.

% ============================================================
% 11 FUTURE WORK
% ============================================================

% \section{\underline{Future Work}}

% This study has several limitations that should be considered when evaluating the results.

% First, the difficulty-prediction experiment was evaluated using a sample of only 100 new problems. Although it gives an initial idea of the model performance  on unseen data, a large sample would allow for more precise and reliable evaluation

% Second, the TF-IDF vectorization approach focuses on the appearance of certain words in the problem statement and fails to recognize deeper logic and the algorithmic background needed to solve the problem \cite{salton1988,manning2008}.

% Third, Random Forest regression technique is not originally designed for high-dimensional texts. In further research, it is necessary to test its performance against both linear methods and transformer-based language models \cite{devlin2019,feng2020,strubell2019}.

% Another limitation relates to the problem tag prediction task. Because a single problem may have multiple tags, and their frequency is very unbalanced (some of them appear much more frequently than others), this makes predicting rare tags more difficult.

% \section{Threats to Validity and Ethical Use}

% There are several limitations that should be considered when interpreting the
% results of this study. First, the original source and collection procedure of
% the dataset are not fully documented. Therefore, it is difficult to make
% strong claims about the completeness and overall coverage of the dataset.

% Another limitation comes from the way the dataset was divided into training
% and testing sets. Since a random problem-level split was used, problems from
% the same Codeforces contest may appear in both sets. This may cause some test
% problems to be more similar to the training problems. In addition, several
% rare algorithmic tags were removed because they did not contain enough
% examples for reliable model training.

% Codeforces ratings should also be interpreted carefully because a rating may
% be influenced by the contest environment, participant performance, and other
% factors. Moreover, the difficulty prediction model was evaluated using only
% 100 unseen problems. Although this gives an initial idea of the model's
% performance, a larger independent test set or cross-validation approach would
% provide a more reliable evaluation of its generalization ability.

% There are also some ethical considerations related to the use of CPVision-AI.
% Predicted algorithmic tags can sometimes provide hints about the possible
% solution approach of a problem. For this reason, the system is mainly intended
% for educational purposes, practice organization, problem analysis, and
% post-contest use rather than during active competitive programming contests.

% Finally, the predictions and confidence scores produced by the system should
% be treated as estimates rather than guaranteed results. The model may produce
% incorrect ratings or tags, especially for problems that are very different
% from those included in the training dataset.
% ============================================================
% 12 CONCLUSION
% ============================================================

\section{\underline{Conclusion}}

This paper presented a lightweight NLP framework for analyzing competitive
programming problems using textual information. The system combines a
unified textual representation, TF--IDF feature extraction, and Random
Forest regression to predict numerical problem difficulty.Using 5,000 training problems and 100 completely unseen test problems from
a dataset of 7,185 Codeforces problems, the difficulty model achieved an
MAE of 112.75 and an RMSE of 251.72. The predicted rating was within
$\pm100$ points of the actual rating for 74.00\% of the test problems and
within $\pm200$ points for 79.00\% of the problems.The project also investigated multi-label algorithmic tag prediction. The
tag prediction component achieved an exact-match accuracy of 1.39\%,
micro-precision of 33.00\%, micro-recall of 50.00\%, and micro-F1 of
39.70\%, nearly doubling the
20.42\% micro-F1 of a majority-tag baseline. These results indicate that
predicting individual algorithmic labels from textual information is
feasible, but recovering the complete set of tags without errors remains
challenging.

Overall, the experiments demonstrate that classical NLP techniques can
serve as a practical and computationally efficient foundation for
competitive programming problem analysis. Future improvements can combine
the current lightweight approach with semantic language models and
structured problem information to improve both difficulty estimation and
algorithmic tag prediction.

% ============================================================
% REFERENCES
% ============================================================

% Allow the references to use the remaining space on the final page.
% A forced page break here made the paper one page longer than necessary.
\balance

\begin{thebibliography}{00}

\bibitem{salton1988}
G. Salton and C. Buckley,
``Term-weighting approaches in automatic text retrieval,''
\emph{Information Processing \& Management},
vol. 24, no. 5, pp. 513--523, 1988,
doi: 10.1016/0306-4573(88)90021-0.

\bibitem{breiman2001}
L. Breiman,
``Random forests,''
\emph{Machine Learning},
vol. 45, pp. 5--32, 2001,
doi: 10.1023/A:1010933404324.

\bibitem{tsoumakas2007}
G. Tsoumakas and I. Katakis,
``Multi-label classification: An overview,''
\emph{International Journal of Data Warehousing and Mining},
vol. 3, no. 3, pp. 1--13, 2007.

\bibitem{codeforcesapi}
Codeforces,
``Codeforces API.''
[Online]. Available:
\url{https://codeforces.com/apiHelp}.
Accessed: Aug. 28, 2026.

\bibitem{kagglecodeforces}
D. I. George,
``Codeforces competitive programming dataset,''
Kaggle, 2025. [Online]. Available:
\url{https://www.kaggle.com/datasets/dinuiongeorge/codeforces-competitive-programming-dataset}.
Accessed: Sep. 5, 2026.

\bibitem{sklearn2011}
F. Pedregosa \emph{et al.},
``Scikit-learn: Machine learning in Python,''
\emph{Journal of Machine Learning Research},
vol. 12, pp. 2825--2830, 2011.

\bibitem{manning2008}
C. D. Manning, P. Raghavan, and H. Sch\"{u}tze,
\emph{Introduction to Information Retrieval}.
Cambridge: Cambridge University Press, 2008.

\bibitem{mccallum1998}
A. McCallum and K. Nigam,
``A comparison of event models for naive bayes text classification,''
in \emph{AAAI-98 Workshop on Learning for Text Categorization},
vol. 752, pp. 41--48, 1998.

\bibitem{devlin2019}
J. Devlin, M.-W. Chang, K. Lee, and K. Toutanova,
``BERT: Pre-training of deep bidirectional transformers for language
understanding,''
in \emph{Proceedings of the 2019 Conference of the North American
Chapter of the Association for Computational Linguistics:
Human Language Technologies},
pp. 4171--4186, 2019.

\bibitem{feng2020}
Z. Feng \emph{et al.},
``CodeBERT: A pre-trained model for programming and natural languages,''
in \emph{Findings of the Association for Computational Linguistics:
EMNLP 2020},
pp. 1536--1547, 2020.

\bibitem{strubell2019}
E. Strubell, A. Ganesh, and A. McCallum,
``Energy and policy considerations for deep learning in NLP,''
in \emph{Proceedings of the 57th Annual Meeting of the
Association for Computational Linguistics},
pp. 3645--3650, 2019.

\bibitem{athavale2019}
V. Athavale, A. Naik, R. Vanjape, and M. Shrivastava,
``Predicting algorithm classes for programming word problems,''
in \emph{Proc. 5th Workshop on Noisy User-generated Text (W-NUT)},
Hong Kong, China, 2019, pp. 84--93,
doi: 10.18653/v1/D19-5511.

\bibitem{iancu2019}
B. Iancu, G. Mazzola, K. Psarakis, and P. Soilis,
``Multi-label classification for automatic tag prediction in the context of
programming challenges,''
\emph{arXiv preprint arXiv:1911.12224}, 2019.

\bibitem{lokat2023}
T. Lokat, D. Prajapati, and S. Labde,
``Tag prediction of competitive programming problems using deep learning
techniques,''
\emph{arXiv preprint arXiv:2308.01863}, 2023.

\bibitem{lobanov2023}
A. Lobanov, E. Bogomolov, Y. Golubev, M. Mirzayanov, and T. Bryksin,
``Predicting tags for programming tasks by combining textual and source code
data,''
\emph{arXiv preprint arXiv:2301.04597}, 2023.

\bibitem{kim2023}
J. Kim, E. Cho, J. Kim, and D. Na,
``Problem-solving guide: Predicting the algorithm tags and difficulty for
competitive programming problems,''
\emph{arXiv preprint arXiv:2310.05791}, 2023.

\bibitem{zhao2018}
W. X. Zhao, W. Zhang, Y. He, X. Xie, and J.-R. Wen,
``Automatically learning topics and difficulty levels of problems in online
judge systems,''
\emph{ACM Transactions on Information Systems}, vol. 36, no. 3,
Art. no. 27, pp. 1--33, 2018,
doi: 10.1145/3158670.

\bibitem{yoshida2024}
C. Yoshida, M. Matsushita, and Y. Higo,
``Estimating the difficulty of programming problems using fine-tuned LLM,''
in \emph{Proc. IEEE/ACIS 22nd International Conference on Software
Engineering Research, Management and Applications (SERA)}, 2024,
pp. 28--34,
doi: 10.1109/SERA61261.2024.10685596.

\bibitem{xu2022}
Y. Xu, Q. Ni, S. Liu, Y. Mi, Y. Yu, and Y. Hao,
``Learning style integrated deep reinforcement learning framework for
programming problem recommendation in online judge system,''
\emph{International Journal of Computational Intelligence Systems},
vol. 15, Art. no. 114, 2022,
doi: 10.1007/s44196-022-00176-4.

\end{thebibliography}
\end{document}
